\SourceLeaf{317}%
Schließlich ergibt sich
\[
\begin{aligned}
\log\lambda&=2.7315487\\
\log l&=2.7315438\\
l&=538''\,9442=0^0\,8'\,58''\,9442.
\end{aligned}
\]

\indent Wie bequem dieses Verfahren ist, merkt man allerdings erst in vollem Umfang, wenn man sich die kleinen Hilfsmittel der Rechenmechanik für solche Verfahren zu eigen gemacht hat. Eine Anleitung dazu wäre hier fehl am Platz. Ich weise nur darauf hin, dass man das, was im obigen Beispiel wie vier Rechnungen aussieht, als eine einzige Rechnung schreiben soll. Bei jeder neuen Überarbeitung ergänzt oder verbessert man dabei nur die letzten Ziffern. Jedenfalls braucht man immer nur die letzte Rechnung aufzubewahren. Darin liegt ein großer Vorteil, besonders bei umfangreichen Vermessungen. Man hat dann für alle Dreiecksseiten den ganzen wesentlichen Kern der Rechnung auf kleinstmöglichem Raum und in der übersichtlichsten Form vor sich, sodass man die Richtigkeit nach Belieben prüfen kann.

\Article{21}
Ich gehe nun zur Hauptaufgabe selbst über. Sie verlangt für die Ellipsoidoberfläche ein ähnliches Verfahren wie das oben für die Kugeloberfläche angegebene. Die Lösung dieser allerdings etwas verwickelten Aufgabe soll hier auf zwei ganz verschiedenen Wegen hergeleitet werden. Die zuerst gegebene Herleitung beruht auf der konformen Abbildung der Ellipsoidoberfläche auf die Kugeloberfläche, deren Theorie in der ersten Abhandlung entwickelt wurde. Die Gewinnung dieser Lösung kann daher als die \textit{erste mittelbare} Anwendung dieser Theorie für die Zwecke der höheren Geodäsie gelten. (Vgl. Artikel 11.)

Es seien nun \(B+\frac12b\) und \(B-\frac12b\) die Breiten zweier Punkte auf der Ellipsoidoberfläche, \(l\) ihr Längenunterschied und \(r\) das zwischen ihnen liegende Stück einer geodätischen Linie (hier in einer beliebigen Einheit gemessen). Die Azimute der Linie am ersten und am zweiten Endpunkt seien \(T+\frac12t\) und \(T-\frac12t\pm180^0\). Es geht also darum, \(b,l\) und \(t\) aus \(r,B\) und \(T\) durch Formeln zu bestimmen, die den oben
